\section{Metodología de investigación aplicada}
El presente proyecto se desarrolló bajo un enfoque cualitativo, orientado a la comprensión e interpretación de conceptos relacionados con la Programación Orientada a Objetos (POO), el patrón Modelo-Vista-Controlador (MVC) y la metodología ágil Scrum. Estos enfoques metodológicos y arquitectónicos ya habían sido definidos desde el inicio del proyecto formativo como parte de la estrategia de desarrollo del sistema de gestión de eventos.

Sin embargo, con el fin de sustentar y validar teóricamente estas elecciones, se procedió a la búsqueda, selección y análisis de artículos científicos que aportaran evidencia sobre sus beneficios y buenas prácticas en el ámbito del desarrollo de software.

\subsection{A. Fase 1- Recolección y construcción del thesaurus}
• Se utilizó como base de datos Google Scholar, realizando consultas con palabras clave relacionadas: patrones de diseño, MVC, POO, Scrum, metodologías ágiles.
• Cada artículo relevante fue documentado en un archivo "Thesaurus" con campos estandarizados: término, fecha, fuente, nombre del artículo.
• El propósito de esta fase fue reunir evidencia científica que respalde el uso de las tecnologías y metodologías seleccionadas en el proyecto formativo.

\subsection{B. Fase 2- Análisis cualitativo de la evidencia (síntesis del thesaurus)}
• El objetivo fue transformar la información documental en argumentos sólidos que justificaran el uso de MVC, POO y Scrum en el desarrollo del sistema, para ello se realizaron lecturas críticas, resúmenes y reflexiones de los artículos, identificando ahí mismo la separación de responsabilidades, mantenibilidad y el trabajo colaborativo en equipos ágiles.
• La evidencia encontrada reforzó la pertinencia de las decisiones iniciales del proyecto y aportó lineamientos prácticos para guiar el diseño e implementación del sistema.